\section{Conclusiones}
\label{sec:Conclusions}


En este trabajo hemos propuesto StreamForest, una arquitectura novedosa para streaming video understanding que aborda las limitaciones en memoria de largo plazo y percepción fina. Al introducir Persistent Event Memory Forest, nuestro método gestiona eficazmente la información visual histórica mediante una fusión adaptativa guiada por penalizaciones de distancia temporal, similitud de contenido y conteo de fusiones. Acoplado con Fine-grained Spatiotemporal Window, el modelo mantiene una comprensión precisa de la escena actual. También presentamos OnlineIT, un dataset de fine-tuning para streaming video understanding que mitiga problemas de desplazamiento espaciotemporal y mejora la percepción y el razonamiento en tiempo real, así como ODV-Bench, un nuevo benchmark adaptado a escenarios de conducción autónoma en tiempo real. Experimentos extensivos demuestran que StreamForest no solo supera a los streaming video MLLMs de estado del arte, sino que también compite con los principales offline video MLLMs bajo configuraciones estrictas de entrada en streaming, mostrando su robustez y valor práctico en aplicaciones de streaming video understanding en tiempo real.
